%% ============================================================
%% VALIDATION PROTOCOL -- split 2 of 3.
%% Drop-in for <<KEEP:stage1>> <<KEEP:stage2>>.
%% Insert BEFORE the existing drafted "validation resolution bounds claimable fidelity"
%% paragraph, which this file does not duplicate and which follows on naturally.
%% Source: v1 lines 1138-1201.
%%
%% NOTE: the gate values quoted here must match tab_val_targets. Per
%% VALIDATION_APPROACH.md the ex-ante gates should be restored to their v1 values
%% (return-temperature MAE 2.5 K, not 1.0 K) before this is merged.
%% ============================================================

The network was upgraded with the heat pump, electrode boiler and thermal storage
\emph{after} the measurement record was collected, so validation follows a two-stage
strategy after \citet{Kus2025} and \citet{Maldonado2024}.

\paragraph{Stage 1: direct network validation.}
The pipe model is validated against pre-upgrade monitoring data with the new assets set to
zero capacity, calibrated by the branch-sequential $U_p$-tuning procedure of
\citet{Maldonado2024} but constrained, in this revision, to a defensible coefficient range
for the pipe class. Table~\ref{tab:valtargets} lists the key performance indicators and
their gates, fixed \emph{ex ante} from published practice before the calibration was run.

How those gates are used differs from the original submission, and the difference is
deliberate. They are reported in full in Section~\ref{subsec:validation}, including the
several that are not met. The temperature gates in particular cannot be met by any model of
this network: the consumer sensors are billing instruments sited downstream of three-way
mixing valves, so the metered temperature sits systematically below the primary junction
temperature, and a point-in-time match is additionally sensitive to a network flow that
billing metering does not pin down. Retaining gates the instrumentation cannot support, and
reporting the failures against them, is the point rather than an oversight -- it is the
evidence for the claim that a model resolved more finely than its metering cannot have that
resolution verified. The conclusions are anchored instead to the annual delivered energy,
and hence the annual network loss, which is gated, met, and the quantity the cost results
depend on.

\paragraph{Stage 2: necessary-condition verification.}
Dispatch of the new assets is verified against plausibility bounds rather than
measurements: heat-pump coefficient of performance within $[2.5, 5.5]$, thermal-storage
state of charge within $[0.05, 0.95]\bar{E}_s$, no simultaneous charge and discharge, and
per-step energy closure. These are necessary conditions, not sufficient ones, and we label
them as such; a post-upgrade measurement campaign is the natural follow-up and is listed
among the limitations.

\paragraph{Structural sanity checks.}
Two checks confirm internal consistency. The nonlinear reference with its quadratic
constraints deactivated reproduces the baseline objective to within $0.01$\,\%; and the
relaxation property $\mathrm{cost}(\CP) \le \mathrm{cost}(\Lone)$ holds on every one of the
135 synthetic instances, since a copperplate can never cost more than the resolved model of
the same network. A violation of either would indicate a formulation error rather than a
finding.
